% Source: https://artofproblemsolving.com/community/c3h2120336_chocolate_eclairs
% Stable identity: aops-topic-2120336-q01
Let Xavier and Yaxley start with $X$ and $Y$ dollars. Then $X+Y=112$ and $X^2=2Y$, so $X^2+2X-224=0$, i.e. $(X-14)(X+16)=0$. Hence $X=14$ and $Y=98$.

Let each buy $p$ pizza slices, and let Yaxley buy $e$ eclairs. Afterwards Xavier has $X'=14-2p$ and Yaxley has $Y'=98-2p-3e$ dollars, both nonnegative integers, and
\[Y'^2=2X'=28-4p .\]
Since $X'\ge0$, $p\in\{0,1,\dots,7\}$, and $28-4p$ takes the values $28,24,20,16,12,8,4,0$. It must be a perfect square, so:
\begin{itemize}
\item $p=3$, $Y'=4$: then $3e=98-6-4=88$, not a multiple of $3$; impossible.
\item $p=6$, $Y'=2$: then $3e=98-12-2=84$, so $e=28$.
\item $p=7$, $Y'=0$: then $3e=98-14-0=84$, so $e=28$.
\end{itemize}
So Yaxley bought $28$ eclairs in every possible case, and the greatest number is $\boxed{28}$.
